% Generated draft: requires mathematical and source review.
\documentclass[11pt]{article}
\usepackage[margin=25mm]{geometry}
\usepackage{fontspec}
\setmainfont{Latin Modern Roman}
\usepackage{amsmath,amssymb,mathtools,unicode-math}
\setmathfont{Latin Modern Math}
\usepackage{enumitem,needspace,xurl,hyperref}
\setlength{\parindent}{0pt}
\setlength{\parskip}{5pt}
\setlength{\emergencystretch}{3em}
\newcommand{\N}{\mathbb{N}}
\newcommand{\Z}{\mathbb{Z}}
\newcommand{\Q}{\mathbb{Q}}
\newcommand{\R}{\mathbb{R}}
\newcommand{\C}{\mathbb{C}}
\newcommand{\F}{\mathbb{F}}
\newcommand{\A}{\mathbb{A}}
\newcommand{\PP}{\mathbb{P}}
\newcommand{\E}{\mathbb{E}}
\newcommand{\eps}{\varepsilon}
\newcommand{\dd}{\,\mathrm{d}}
\begin{document}
\section*{UnsolvedMath ID: 7200081}
\textbf{Source problem code:} catalogue entry
\subsection{Definitions and mathematical statement}
For an embedded closed $C^{1,1}$ curve $\gamma:S^1\to\R^3$, let $L(\gamma)$ be its length. In the tube-radius convention, its thickness $\tau(\gamma)$ is its reach: the supremum of radii $r$ such that every point at distance less than $r$ from the image has a unique nearest point on the curve. For positive thickness define
\[
\operatorname{Rop}(\gamma)=\frac{L(\gamma)}{\tau(\gamma)}.
\]
This ratio is invariant under Euclidean scaling. For a tame knot type $K$, the minimum-ropelength invariant is conventionally
\[
\operatorname{Rop}(K)=\inf\{\operatorname{Rop}(\gamma):\gamma\text{ is an embedded closed curve of knot type }K\}.
\]
The admissible regularity class and attainment convention must be fixed; the infimum notation does not by itself prove attainment.

The source asks for a general expression for this value for an arbitrary closed knot. It does not define what qualifies as an expression: an exact formula from a knot diagram or finite invariant, a terminating exact algorithm, and arbitrarily accurate numerical approximation are distinct targets. Merely restating the variational infimum is not a proposed solution. The thickness convention is also unstated; defining rope thickness by diameter rather than radius rescales all numerical values.

The supplied summary claims exact values for the trefoil and figure-eight without any proof or reference. Numerical tight-knot candidates and rigorous global minima are different information, and those unsupported exact-value assertions are not adopted. The central request remains a formula/classification programme with its output and computability criteria unspecified.
\subsection{Short English statement}
Find a general way to express the least length-to-thickness ratio in each closed knot type. Specify the thickness convention and whether the desired output is an exact formula, exact algorithm, or approximation.
\subsection{Sources}
\begin{itemize}
\item Local UnsolvedMath registry, permanent ID 7200081, exact source code catalogue entry; no proposer or year supplied. Registry status open, classification OPEN-TRIAGE.
\item Wikipedia contributors, \emph{List of unsolved problems in mathematics}, Geometry bullet 81, pinned revision 1366636767, \url{https://en.wikipedia.org/w/index.php?title=List_of_unsolved_problems_in_mathematics&oldid=1366636767#Geometry}. Supplied MediaWiki extraction accessed July 29, 2026. No primary citation for the general-expression question is given.
\item The summary cites ``The ropelength conjecture of alternating knots'', \texttt{arXiv:2208.00123}, as related crossing-number work, without authors, and reports only abstract-level inspection. The identifier/title pairing and theorem scope are unverified here. It reports no 2024--2026 general closed form, but its additional exact-trefoil/figure-eight claims have no supporting citation or optimality certificate.
\item Provenance flags primary-source status review NEEDS\_REVIEW and says the bullet remains unpublished pending review despite published=true metadata. Wikipedia text is CC BY-SA 4.0; attribution and share-alike handling remain required for release.
\end{itemize}
\end{document}
